\documentclass{article}
\usepackage{graphicx} % Required for inserting images
\usepackage{amssymb}
\usepackage{bm}
\usepackage{bbm}
\usepackage{hyperref}
\usepackage{amsmath}
\usepackage{algorithm}
\usepackage{algpseudocode}
\usepackage{tikz-cd}
\usepackage{float}


\title{Next token prediction}
\author{Jake Feldman Starosta}
\date{June 2026}

\begin{document}

\maketitle

\section{Introduction}
We want to model feature-label pairs $(\bm{x}^k, y^k)$ in a dataset $\mathcal{D}$, where $y^k$ is the empirical conditional probability measure given that features are $N$th-order Markov. 

\subsection{Notation:}
\begin{itemize}
\item $S \subset \mathbb{R}^d$ - individual token embedding space 
\item $\mathcal{N} = \{1, ..., N\} $ - token index of $N$th-order Markov chain
\item $x^{i,k} \in S$ - individual token embedding
\item $\tilde y^k \in S$ - target token embedding
\item $\bm{x}^k = (x^{1, k}, \dots, x^{N, k}) \in S^N$ - training sample/feature vector
\item $p_i = \frac{i}{N} \in PE_N$ - positional encoding of token $i$
\item $X^{i,k} = (x^{i, k} , p_i) \in \mathbb{X} = S\times PE_N$ - feature with token and  positional encoding
\item $U = (W, A, b, Q, K, V) \in \mathbb{U}$ - actions/model weights


\end{itemize}

The difference between this problem and the original paper's \cite{akman2026optimalcontrolapproachtransformer} formulation is that our final distribution has 1. no positional information and 2. lives in a different space to our input. 
\section{Dataset}
We begin by constructing our desired dataset with probability measure labels based on a given dataset with deterministic token embedding labels resembling those of the original paper (however, we consider singletons, not sequences).

An initial dataset is given as 
$$\tilde D = \{(\bm{x}^k, \tilde y^k):k \in \tilde {\mathcal{K}} = \{1, \dots , \tilde K\}\} \subset S^N \times S$$
where $\bm{x}^k$ is a feature vector and $\tilde{y}^k \in S$ is the target token embedding. The initial dataset $\hat D$ has training samples indexed by $\tilde {\mathcal{K}}$, with possibly some duplicate feature vectors, i.e., $\bm{x}^j = \bm{x}^l$ for $j\neq l$, and $j, l \in \tilde {\mathcal{K}}$. We now construct a new dataset $D$ from $\tilde D$ with labels as probability measures instead of individual token embeddings. 
We index the unique feature vectors with $\mathcal{K} = \{1, \dots, K\}$, with one representative index picked from $\tilde {\mathcal{K}}$ for each unique feature vector. For $k \in \mathcal{K}$, we define the empirical measure $y^k \in P(S)$ conditioned on $\bm{x}^k$ as  
$$y^k = \frac{
\sum_{l=1}^{\tilde {K}} \delta _ {\tilde y^l} \mathbbm{1}_{\{\bm{x}^l = \bm{x}^k\}}
}{\sum_{l=1}^{\tilde {K}} \mathbbm{1}_{\{\bm{x}^l = \bm{x}^k\}}
}
.$$

Then $D = \{(\bm{x}^k, y^k):k \in \mathcal{K}\} \subset S^N \times P(S)$ has unique training samples $\bm{x}^k$, with labels $y^k$ as empirical probability measures. Note that the denominator is non-zero because it is the count of samples in $\hat D$ with feature vector $\bm{x}^k$.

An important distinction between this dataset construction and that of large language models is that here we identify a single target distribution label for a given sample of $N$ tokens. Ideally, we would like to use the information from each token to train our next token prediction model: for an $N$ word sample, we would create $N-1$ sample-label pairs by masking the tokens at and after the predicted token.

\section{Transformer Dynamics}
Per the paper, $\bm{x}^k = (x^{1, k}, \dots, x^{N, k})$ gives rise to $X^k = (X^{1, k}, \dots X^{N, k})$ where $X^{i, k} = ( p_i, x^{i, k})$ for $i \in \mathcal{N} := \{1, \dots, N\}$. Again, we construct an empirical measure $\mu^i_t \in P(\mathbb{X})$, encoding each particle state and position information. It evolves with a deterministic pushforward $\Phi: P(\mathbb{X})\times \mathbb U \to P(\mathbb{X})$ for $T$ layers indexed by $\mathcal{T}$. We then aggregate all $\mu^i_t$ to an ensemble $((\mu^1_t, ... , \mu^K_t))_{t \in \mathcal{T}}=(\vec \mu_t)_{t \in \mathcal{T}}$ with transition kernel $\eta : P(\mathbb{X})^K \times \mathbb U\to P(P(\mathbb{X}))^K$, where $\eta (\cdot|\vec\mu, U) = \delta _{(\bm{\Phi} (\vec \mu, U))}(\cdot)$ and $\bm{\Phi} : P(\mathbb{X})^K\times \mathbb U \to P(\mathbb{X})^K$ with $\Phi$ applied pointwise to each $\mu^i_t$. 

This defines a MDP on the state space $P(\mathbb{X})^K$ having elements $(\vec \mu_t)_{t \in \mathcal{T}}$, with action space $\mathbb U$, transition kernel $\eta$ and the cost $c$ shown below.

We aim to minimize the following loss functional for a policy $\gamma$:

$$
V(D, \gamma) =  \frac{1}{K}\sum _{k = 1}^K W_{2}(\nu^k ,y^k)^2 
$$

We define $\nu^k := \delta_{x^{k, N}_{T, \gamma}}$, where $x^{k, N}_{T, \gamma}$ starts at $x^{k, N}_{0, \gamma} \in S$ belonging to the ensemble probability measure $\mu^k$ evolving according to $\bm{\Phi}$ which takes actions from $\gamma$. We have

$$W_2(P, Q) = \inf _{\pi \in \Pi (P, Q)} \int _{S \times S} \pi(dx, dy) c(x, y)^2$$

with the cost function 

$$ c(x, y) =  \sqrt{||x-y||_2^2}.$$

For our loss, we only compare the probability measure of the last position ($p_N = \frac{N}{N} = 1$) in the output vectors to our label. This is also a trick used in practice in transformer fine-tuning, where features and labels are separated this way (see equation 3 of \cite{radford2018improving}). 

Although this works for training, we're not actually predicting a distribution. We're just predicting a token embedding, and comparing it to the empirical conditional probability measure. 

\section{Quantization}

We select parameters $l, m, n$ to quantize the above MDP into MDP$^{(l, m, n)}$ as follows: we quantize the values of the probability measure on the state space with $l$, the state space itself with $n$, and the action space with $m$. 
We quantize the token embedding space as $S_n$ such that given $x \in S$, we can find some $s \in S_n$ that satisfies
$$|x - s| < \frac{1}{n}$$
which, due to the compactness of $S$, exists and is finite for all $n$. Similarly, due to the compactness of $\mathbb{U}$, we quantize the action space finitely as $\mathbb{U}_m$ so that given $U \in \mathbb{U}$,
$$||U - \hat U||_F < \frac{1}{m}$$
for some $\hat U \in \mathbb{U}_m$, where $||\cdot ||_F$ is the product Frobenius norm. 

Let $\mathbb{X}_n := PE_N \times S_n$ and define the set of quantized probability measures on $\mathbb{X}_n$ as 
$$\mathcal{P} ^{(l)}(\mathbb X_n) = \{(p_a)_{a = 1}^{|\mathbb X_n|}: p_a \in \frac{1}{l} \mathbb{N}_0, \sum_{a = 1}^{|\mathbb X_n|} p_a = 1\}.$$

We define the quantizers $Q_n : \mathbb{X} \to \mathbb{X}_n$ mapping $x \in \mathbb{X}$ to the closest point in $\mathbb{X}_n$ while preserving position, and $R_l : \mathcal P(\mathbb X _n) \to \mathcal P^{(l)}(\mathbb X _n)$ mapping $\hat \mu \in \mathcal P(\mathbb X _n)$ to the closest element in $\mathcal P^{(l)}(\mathbb X _n)$. 

The quantized deterministic ensemble pushforward $\bm \Phi^{(l, n, m)}$ maps $\mathcal P(\mathbb{X}_n)^K\times \mathbb U_m \to \mathcal P(\mathbb X_n)^K$ by quantizing and applying the single-measure pushforward $\Phi ^{(n)}: \mathcal P(\mathbb X_n) \times \mathbb U\to \mathcal P (\mathbb X_n)$ point-wise. It induces the transition kernel $\eta ^{(l, n, m)} : \mathcal P(\mathbb{X}_n)\times \mathbb U_m \to \mathcal P(\mathcal P(\mathbb X_n))$. This characterizes $\text{MDP}^{(l, n, m)}$ with a state space $\mathcal P(\mathbb{X}_n)^K$, action space $\mathbb U_m$, transition kernel $\eta^{(l, n, m)}$ and the same cost $c$ as MDP.

For $\hat \mu, \hat \nu \in P(S_n)$, we define the measure-level cost as 
$$\hat W_2 (\hat \mu, \hat \nu) = \min _{\hat \pi \in \Pi (\hat \mu, \hat \nu)} \sum _ {\hat x \in S_n}  \sum _ {\hat y \in S_n} \hat \pi (\hat x, \hat y)c(\hat x, \hat y)^2$$

where $$\hat \Pi (\hat \mu, \hat \nu) = \{\hat \pi \in S_n \times S_n : \hat \pi (a,b) \geq 0, \sum _ {a \in S_n} \hat \pi (a, b) = \hat \nu, \sum _ {b \in S_n} \hat \pi (a, b) = \hat \mu\}.$$ 

Our terminal loss functional for $\text{MDP}^{(l, n, m)}$ is given as 
$$\hat V(D, \hat \gamma) = \frac{1}{K}\sum _{k = 1}^K \hat W_{2}(\hat \mu^{k, N}_T ,\hat y^k)^2$$

where $\hat \mu^{k, N}_T := \delta_{\hat x^{k, N}_{T, \gamma}}$ and $\hat y^k := R_l(y^k)$. Our prediction $\hat x^{k, N}_{T, \hat \gamma}$ is initialized at $\hat x^{k, N}_{0, \hat \gamma} \in S_n$ belonging to the ensemble probability measure $\hat \mu^k$ evolving according to $\bm{\Phi}^{(l, n, m)}$ which takes actions from $\hat \gamma$. 

Note that this loss functional is sufficient for training (as it uses all our prediction information), but does not necessarily translate to performance. Specifically, taking a small $l$ may artificially decrease loss without gaining performance due to coarseness. 

For a given dataset $D$, we pose the dynamic programming equations for $\text{MDP}^{(l, n, m)}$ as follows
$$
\begin{aligned} & \mathfrak C _ {l, n, m, T} (\hat {\bm \mu}) := \frac{1}{K}\sum _{k = 1}^K \hat W_{2}(\hat \mu^{k, N}_T ,\hat y^k)^2\\
&\mathfrak C _ {l, n, m, t} (\hat {\bm \mu}) := \min_{U \in \mathbb U_m} \sum _{\hat{\bm \mu}_{t +1} \in P^{(l)}(\mathbb X_n)^K} \eta^{(l, n, m)}(\hat{\bm \mu}_{t + 1}| \hat{\bm \mu}_t, U)\mathfrak C _ {l, n, m, t + 1} (\hat{\bm \mu}_{t + 1})\\
&\phantom{\mathfrak C _ {l, n, m, t} (\hat {\bm \mu})}= \min_{U \in \mathbb U_m} \mathfrak C _ {l, n, m, t+1} (\bm \Phi ^{(l, n, m)}(\hat {\bm \mu}, U))
\end{aligned}
$$

for all $\hat{\bm \mu} = (\hat \mu ^1, \dots \hat \mu ^K)\in \mathcal P^{(l)}(\mathbb X _n)^K$ and $t \in \mathcal T$. The last equality is due to the deterministic nature of $\bm \Phi ^{(l, n, m)}$.

Note for myself: do we not need to discretize the second one? Not that it should change the policy, but that we don't even have a clean space defined. Also, because the domain of integration isn't continuous. 

\begin{figure}
    \centering
    \begin{tikzcd}[row sep=3em, column sep=3em]
    P(\mathbb{X}) 
      & P(\mathbb{X}_n) \arrow[r, "R_l"] 
      & P^{(l)}(\mathbb{X}_n) 
      \\
    \mathbb{X} \arrow[u, "\text{lift as } \text{probability measure}"] \arrow[r, "Q_n"]
      & \mathbb{X}_n \arrow[u] 
      & 
      \\
    S \arrow[u, "\text{lift with } P_{EN}"] \arrow[r, "\text{quantize}"] 
      & S_n \arrow[u]
      & \\
    \end{tikzcd}
    \caption{Two-stage State Quantization}
    \label{fig:quantization}
\end{figure}

\newpage
\section{Algorithm}

\begin{algorithm}
\caption{Next Token Prediction with Dynamic Programming}
\begin{algorithmic}[1]

\State \textbf{Input:} data set $\tilde D = \{(\bm x ^k, \tilde y^k)\}_{k=1}^K$, horizon $T$, quantization levels $(l, n, m)$

\State Construct state grid $S_n$ from $S$ and for $(p,x) \in \mathbb X$ define 
$$Q_n(p, x) = (p, \arg \min_{s \in S_n}||x-s||)$$

\State Construct finite action set $\mathbb U_m = \{U^1, \dots U^m\}$

\State Construct measure grid $\mathcal P^{(l)} (\mathbb X_n)$ and for $\nu \in P(\mathbb X_n)$ define 
$$R_l(\nu)  = \arg \min _{\hat \nu \in \mathcal P^{(l)} (\mathbb X_n)} W_2(\nu, \hat \nu)$$

\State \textbf{for} $k = 1, \dots , K$ \textbf{do}

\State \quad $\hat \mu_0^k \leftarrow R_l (\frac{1}{N}\sum _{i=1}^N\delta_{Q_n(p_i, x^{i,k})})$

\State \quad $\hat y^k \leftarrow R_l( {\frac{
\sum_{j=1}^{\tilde {K}} \delta _ {\tilde y^j} \mathbbm{1}_{\{\bm{x}^j = \bm{x}^k\}}}{\sum_{j=1}^{\tilde {K}} \mathbbm{1}_{\{\bm{x}^j = \bm{x}^k\}}}})$

\State \textbf{end for}

\State Construct dataset $D = \{(\bm x ^k, \hat y^k)\}_{k=1}^K$

\State Define dynamics for $U \in \mathbb U _m$
$$\Phi ^{(l, n, m)}(\hat \mu, U) = R_l(\Phi^{(n)}(\hat \mu, U))$$

\State Define terminal cost $$\mathfrak C _ {l, n, m, T} (\hat {\bm \mu}) = \frac{1}{K}\sum _{k = 1}^K \hat W_{2}(\hat \mu^{k, N}_T ,\hat y^k)^2$$

\State \textbf{for} $t = T-1$ down to $0$ \textbf{do}
\State \quad \textbf{for all} $\hat{\bm \mu} = (\hat \mu ^1, \dots , \hat \mu ^K) \in \mathcal P^l(\mathbb X_n)^K$ \textbf{do} 

\State \quad \quad 
$$\mathfrak C _ {l, n, m, t} (\hat {\bm \mu}) = \min_{U \in \mathbb U_m} \mathfrak C _ {l, n, m, t+1} (\Phi ^{(l, n, m)}(\hat {\mu}^1, U), \dots ,\Phi ^{(l, n, m)}(\hat {\mu}^K, U))$$

\State \quad \quad 
$$\gamma ^*_t(\hat {\bm \mu}) = \arg \min _{U \in \mathbb{U}_m} \mathfrak C _ {l, n, m, t+1} (\Phi ^{(l, n, m)}(\hat {\mu}^1, U), \dots ,\Phi ^{(l, n, m)}(\hat {\mu}^K, U))$$

\State \quad \textbf{end for}
\State \textbf{end for}
\State \textbf{Forward Pass:}

\State \textbf{for} $t = 0, \dots , T-1$ \textbf{do}

\State \quad $U^*_t \leftarrow \gamma ^*_t(\hat \mu _t ^1, \dots , \hat \mu _t ^K)$

\State \quad \textbf{for} $k = 1, \dots K$ \textbf{do}

\State \quad \quad $\hat \mu^k_{t+1} \leftarrow \Phi ^{(l, n, m)}(\hat \mu^k_t, U^*_t)$

\State \quad \textbf{end for}
\State \textbf{end for}
\State \textbf{Output:} $(U^*_t) ^{T -1}_{t = 0}$

\end{algorithmic}
\end{algorithm}


\section{Next steps}

Here are some possible options:

\begin{itemize}
\item Vocabulary (predefined list to map onto and out of $S$, include unembedding matrix in actions)
\item Masking and sequence-to-sequence prediction (as is actually done in practice)
\item Filter formulation (Semantic hidden state. Token 'path' as the visible approximations of semantics (?))
\item   Determine if the underlying state is learned using next token prediction (?)


\end{itemize}

\bibliographystyle{plain}
\bibliography{References}

\end{document}
